%; whizzy paragraph -pdf xpdf -latex ./whizzypdfptex.sh
%; whizzy-paragraph "^\\\\begin{frame}\\|\\\\emtext"
% latex beamer presentation.
% platex, latex-beamer でコンパイルすることを想定。 

%     Tokyo Debian Meeting resources
%     Copyright (C) 2012 Junichi Uekawa

%     This program is free software; you can redistribute it and/or modify
%     it under the terms of the GNU General Public License as published by
%     the Free Software Foundation; either version 2 of the License, or
%     (at your option) any later version.

%     This program is distributed in the hope that it will be useful,
%     but WITHOUT ANY WARRANTY; without even the implied warreanty of
%     MERCHANTABILITY or FITNESS FOR A PARTICULAR PURPOSE.  See the
%     GNU General Public License for more details.
%     You should have received a copy of the GNU General Public License
%     along with this program; if not, write to the Free Software
%     Foundation, Inc., 51 Franklin St, Fifth Floor, Boston, MA  02110-1301 USA

\documentclass[cjk,dvipdfmx,12pt]{beamer}
\usetheme{Tokyo}
\usepackage{monthlypresentation}

%  preview (shell-command (concat "evince " (replace-regexp-in-string "tex$" "pdf"(buffer-file-name)) "&")) 
%  presentation (shell-command (concat "xpdf -fullscreen " (replace-regexp-in-string "tex$" "pdf"(buffer-file-name)) "&"))
%  presentation (shell-command (concat "evince " (replace-regexp-in-string "tex$" "pdf"(buffer-file-name)) "&"))

%http://www.naney.org/diki/dk/hyperref.html
%日本語EUC系環境の時
\AtBeginDvi{\special{pdf:tounicode EUC-UCS2}}
%シフトJIS系環境の時
%\AtBeginDvi{\special{pdf:tounicode 90ms-RKSJ-UCS2}}

\newenvironment{commandlinesmall}%
{\VerbatimEnvironment
  \begin{Sbox}\begin{minipage}{1.0\hsize}\begin{fontsize}{8}{8} \begin{BVerbatim}}%
{\end{BVerbatim}\end{fontsize}\end{minipage}\end{Sbox}
  \setlength{\fboxsep}{8pt}
% start on a new paragraph

\vspace{6pt}% skip before
\fcolorbox{dancerdarkblue}{dancerlightblue}{\TheSbox}

\vspace{6pt}% skip after
}
%end of commandlinesmall

\title{東京エリアDebian勉強会}
\subtitle{第141回 2016年7月度}
\author{杉本 典充}
\date{2016年7月16日}
\logo{\includegraphics[width=8cm]{image200607/openlogo-light.eps}}

\begin{document}

\begin{frame}
\titlepage{}
\end{frame}

\begin{frame}{Agenda}
 \begin{minipage}[t]{0.45\hsize}
  \begin{itemize}
  \item 最近あったDebian関連のイベント報告
	\begin{itemize}
	\item 第140回東京エリアDebian勉強会
	\end{itemize}
  \item 事前課題発表
  \end{itemize}
 \end{minipage}
 \begin{minipage}[t]{0.45\hsize}
  \begin{itemize}
   \item Debian Trivia Quiz
   \item Rogerさん「Debian Installer Screen対応」
   \item 杉本さん「Debianでlxcをセットアップしてみよう」
  \end{itemize}
 \end{minipage}
\end{frame}

\section{イベント報告}
\emtext{イベント報告}

\begin{frame}{第140回東京エリアDebian勉強会 }

\begin{itemize}
\item 場所はdotsさんをお借りして開催しました
\item 参加者は9名
\item セミナ内容は、以下内容を発表しました
  \begin{enumerate}
  \item kenhysさん「dexexpoをハックするには」
  \end{enumerate}
\item 残りの時間でhack timeを行い、成果発表をしました
\item \url{https://tokyodebianmeeting.titanpad.com/ep/pad/export/140/rev.858?format=txt}
\end{itemize} 
\end{frame}

\section{事前課題}
\emtext{事前課題}
{\footnotesize
 \begin{prework}{ koedoyoshida }
  \begin{enumerate}
  \item $B$"$j$^$9(B
  \item Rpi
  \item $BM=Dj$OL$Dj(B
  \end{enumerate}
\end{prework}

\begin{prework}{ Ko Kazaana }
  \begin{enumerate}
  \item $B$"$j$^$;$s(B
  \item ($B2sEz$J$7(B)
  \item Debian$BNr;KG/I=$N@):n(B
  \end{enumerate}
\end{prework}

\begin{prework}{ NOKUBI Takatsugu }
  \begin{enumerate}
  \item $B$"$j$^$9(B
  \item FUJITSU Server PRIMERGY TX1310 M1
  \item deb$B%Q%C%1!<%8$N(Blintian$B7Y9pBP1~(B
  \end{enumerate}
\end{prework}

\begin{prework}{ yutannihilation }
  \begin{enumerate}
  \item $B$"$j$^$;$s(B
  \item ($B2sEz$J$7(B)
  \item Scaleway$B$G<+:n(BDebian$B%$%a!<%8$r;H$&J}K!$rD4$Y$^$9(B
  \end{enumerate}
\end{prework}

\begin{prework}{ kenhys }
  \begin{enumerate}
  \item $B$"$j$^$9(B
  \item ARM$B%\!<%I(B
  \item $B%Q%C%1!<%8$N%a%s%F%J%s%9$J$I(B
  \end{enumerate}
\end{prework}

\begin{prework}{ dictoss }
  \begin{enumerate}
  \item $B$"$j$^$9(B
  \item $B9uH"(BHG$B!"(Braspberry pi 2B+$B!"(BCubox
  \item preseed$B$rD4$Y$k!"(Bkfreebsd$B4XO"D4::(B
  \end{enumerate}
\end{prework}

\begin{prework}{ Roger Shimizu }
  \begin{enumerate}
  \item $B$"$j$^$9(B
  \item ($B2sEz$J$7(B)
  \item PGP key $B%5%$%s!)(B
  \end{enumerate}
\end{prework}

\begin{prework}{ ymmt2005 }
  \begin{enumerate}
  \item $B$"$j$^$;$s(B
  \item ($B2sEz$J$7(B)
  \item go-apt-cacher/mirror $B$K$D$$$F2?$+$r(B
  \end{enumerate}
\end{prework}

\begin{prework}{ Kouhei Maeda }
  \begin{enumerate}
  \item $B$"$j$^$;$s(B
  \item OpenBlockS266, 600, AX
  \item $B%Q%C%1!<%8%a%s%F%J%s%9(B
  \end{enumerate}
\end{prework}

\begin{prework}{ Charles Plessy }
  \begin{enumerate}
  \item $B$"$j$^$9(B
  \item https://wiki.debian.org/Cloud/AmazonEC2DebianInstaller
  \item $B:#2s$b!"(BGPG$B80$H(BSMTP$B%;%C%H%"%C%W(B
  \end{enumerate}
\end{prework}

}

\subsection{問題}

%; whizzy-master ../debianmeetingresume201311.tex
% $B0J>e$N@_Dj$r$7$F$$$k$?$a!"$3$N%U%!%$%k$G(B M-x whizzytex $B$9$k$H!"(Bwhizzytex$B$,MxMQ$G$-$^$9!#(B
%

\santaku
{quiz1}
{a}
{b}
{c}
{C}
{desc}


\section{Debian Installer Screen対応}
\emtext{Debian Installer Screen対応}

\section{Debianでlxcをセットアップしてみよう}
\emtext{Debianでlxcをセットアップしてみよう}

\section{Hack time}
\emtext{Hack time}

\begin{frame}{成果を記入下さい!}

  今回、Hack time時の成果を記録に残してみます。

\url{https://tokyodebianmeeting.titanpad.com/141}

に皆さんアクセス頂き（認証不要です）、各自成果を18:30までに記録するようにお願いします。

\end{frame}
  
\section{今後のイベント}
\emtext{今後のイベント}
\begin{frame}{今後のイベント}
\begin{itemize}
\item 8/某日 東京エリアDebian勉強会
  \begin{itemize}
  \item 発表者募集中
  \end{itemize}
\end{itemize}
\end{frame}

\end{document}

;;; Local Variables: ***
;;; outline-regexp: "\\([ 	]*\\\\\\(documentstyle\\|documentclass\\|emtext\\|section\\|begin{frame}\\)\\*?[ 	]*[[{]\\|[]+\\)" ***
;;; End: ***
